% page 134 of the facsimile (image p-133 of the scan)
% block 157.5 x 246.3 mm ; left margin 8.0 mm
\pgc{-12.700mm}{0.000mm}{
\l{\marge[78.737mm]{18.471mm}{\ornamento{11.765mm}{ornaments/letters/letter-E.png}}\cel{2}1244}
\l{}
\l{}
\l{}
\l{}
\l{}
\l{}
\l{e. (muziko).Noto sepesma di la \textquotedbl{}do\textquotedbl{}-gamo..}
\l{}
\l{-e-.\cel{2}Infixo qua signifikas \textquotedbl{}samakolora kam\textquotedbl{} o \textquotedbl{}sam-aspektaa}
\l{\cel{4}kam\textquotedbl{}.}
\l{}
\l{\cel{1}-e .\cel{2}Sufixo di la adverbi derivita.}
\l{}
\l{\cel{1}\sou{e(d)}. Kunjunciono qua uzesas por interligar la frazi o}
\l{\cel{4}la propozicioni koordinita.}
\l{}
\l{\cel{1}\sou{ebeno}. (\sou{bot.}) Ligno de la ebeniero (L.\cel{1}\sou{dicapyros ebenum}),}
\l{\cel{4}tre pezoza, tre harda, di qua la grano esas tre densa,}
\l{\cel{4}e remarkinda pro la beleso di lua obskur-nigreso. - DEFS.}
\l{}
\l{-ebl-.\cel{2}Sufixo qua signifikas \textquotedbl{}quan onu povas...\textquotedbl{}..}
\l{}
\l{\cel{1}\sou{ebonito}. Kauchuko hardigita per vulkanizo, de qua onu facas}
\l{\cel{3}pektili, juveli, izolivi por la aparati elektrala.- DEFIRS.}
\l{}
\l{\cel{1}\sou{ebria}. I.Di qua la cerebro-funciono trublesas da drinkir}
\l{\cel{4}tro multe vino od altra drinkajo fermentacita. - II.}
\l{\cel{4}Di qua la spirito trublesas da la deliro de pasiono.-}
\l{\cel{4}EFIS.}
\l{}
\l{ebuliar. (netrans.)Agitesar (aludante la surfaco di liqui--}
\l{\cel{4}do) kande lu bolieskas. - EFIS.}
\l{}
\l{ebulioskopo. Aparato per qua onu determinas, per bolio,,}
\l{\cel{4}la quanto de alkoholo quan kontenas alkoholajo.}
\l{}
\l{\cel{1}\sou{ecelar}. (\sou{netrans}.) (\sou{en, pri}). Esar, en sua genero, ye grado}
\l{\cel{4}eminenta. - DEFI.}
\l{}
\l{\cel{1}\sou{ecelenco}. Titulo honorala, donata a la ministri, a la am-}
\l{\cel{4}basadisti. - DEFIS.}
\l{}
\l{ccentriaa. I. (\sou{geom.}) Di qua la centro eskartesas de pn +n}
\l{\cel{3}donita. - II. Ecentriko : disko uzata por transformarr}
\l{\cel{4}movo cirklo-kurva, per ranuro kalkulita segun la movo}
\l{\cel{4}produktenda rektalinea (iro e retroveno), ed en qua cir-}
\l{\cel{4}kulas roteto qua komandas la mekanismo. - DEFIS.}
\l{}
\l{\cel{1}\sou{eceptar}. (\sou{trans}.) Lasar extere persono, objekto, pri qua}
\l{\cel{4}onu afirmas ulo. - EFIS.}
\l{}
\l{\cel{1}\sou{ecesar}. (\sou{trans}.) I. Transirar limito fixigita.\cel{2}II. Trans-}
\l{\cel{4}irar to quon ulu povas tolerar. - DEFIRS.}
\l{}
\l{\cel{1}\sou{echekar}. (\sou{trans}.) Kontreagar vinkoze sua adverso, e min-}
\l{\cel{4}acar lu.}
}
